\documentclass{article}
\usepackage{graphicx}
\usepackage{amsmath}
\usepackage{amssymb}
\usepackage{amsfonts}
\usepackage{hyperref}
\usepackage{exercises}
\usepackage{indentfirst}
\usepackage{circledsteps}
\usepackage{amsthm}
\usepackage{csquotes}
\usepackage{yhmath}
\usepackage{comment}
\usepackage{placeins}
\usepackage{geometry}
\usepackage{hyperref}
\usepackage{tikz}
\usepackage{tabto}
\newcommand*{\oldneg}{\mathord{\sim}}
\newcommand{\Z}{\mathbb{Z}}
\newcommand{\N}{\mathbb{N}}
\let\oldemptyset\emptyset

\theoremstyle{definition}
\newtheorem{definition}{Definition}
\newtheorem{theorem}{Theorem}[section]
\newtheorem{corollary}{Corollary}[theorem]
\newtheorem{lemma}[theorem]{Lemma}

\setlength{\parindent}{.5in}
\geometry{margin=0.5in}
\setlength{\parskip}{12pt}
\title{Chapter 15 Notes}
\begin{document}
\date{}
\maketitle

\section*{Section 0}

\subsection*{Exercise 0.1}
I will note that I was not able to prove this unless we assume that $A$ has some cardinality (or rather, that $A$ has a choice function). The 
parts (a)-(c) of this exercise basically follows notes I had in the \textquote{Theorem 1.1} of Section 1 of the Chapter 11 notes.

(d) We show by transfinite induction on all infinite ordinals that for each infinite ordinal $\alpha$, the set $L_\alpha$ has 
cardinality at most $\overline{\overline{\alpha}}$.

For the successor case, suppose the result holds for some $L_\alpha$. If $\alpha+1$ is infinite, then $\alpha$ must also be infinite. By part (c) 
of the exercise, it follows that $\overline{\overline{L_{\alpha+1}}}=\overline{\overline{L_\alpha}}$, which completes the proof in this case.

For the limit case, suppose that for all ordinals $\beta$ less than some limit ordinal $\alpha$, $\overline{\overline{L_\beta}}\le\beta$.
Because $\alpha$ is a limit ordinal, $L_\alpha=\cup_{\beta<\alpha}L_\beta$. If $L_\alpha=L_\omega$, we are done, as $L_\omega$ is denumerable.
Otherwise, it is easy to show that all $L_\beta$ for $\beta<\alpha$ are such that $\overline{\overline{L_\beta}}\le\overline{\overline{\alpha}}$.
Because $L_\alpha=\cup S$ where $S$ is the set of all $L_\beta$ with $\beta<\alpha$, it follows that $\overline{\overline{L_\alpha}}\le\overline{\overline{\alpha}}$ 
(because each memeber of $S$ has at most cardinality $\overline{\overline{\alpha}}$, and there are $\alpha$ members of $S$).

(e)
For this problem, note that if $\alpha$ is an infinite ordinal, the number of successor ordinals in $\alpha$ must be 
$\overline{\overline{\alpha}}$ (Reason: It cannot be that the number of successor ordinals in $\alpha$ was less than $\overline{\overline{\alpha}}$ 
and the number of limit ordinals in $\alpha$ was less than $\overline{\overline{\alpha}}$, or else the number of elements in 
$\alpha$ would be less than $\overline{\overline{\alpha}}$. If the number of limit ordinals in $\alpha$ was $\overline{\overline{\alpha}}$, then there are two cases: Either 
$\alpha$ is the successor of a limit ordinal, or it isn't. If it isn't, then for every limit ordinal $\lambda$ in $\alpha$, $\lambda^+$ is also in $\alpha$, meaning 
there are also $\overline{\overline{\alpha}}$ successor ordinals in $\alpha$. Otherwise, if $\alpha$ is the successor of a limit ordinal, there still are 
$\overline{\overline{\alpha}}$ successor ordinals in $\alpha$, as we could just take any successor ordinal less than $\omega$, which replaces the successor ordinal 
after $\lambda$ for $\lambda$ the predecessor of $\alpha$. Thus, if there are $\overline{\overline{\alpha}}$ limit ordinals in $\alpha$, there are also
$\overline{\overline{\alpha}}$ successor ordinals in $\alpha$. Hence, in any case, there are $\overline{\overline{\alpha}}$ successor ordinals in $\alpha$). 
Now consider some infinite cardinal $c$. Then $c$ contains $c$ successor ordinals. Then, for each 
successor ordinal $\alpha^+<c$, $L_\alpha$ contains some element not in $L_\alpha$ (for example, $L_\alpha$ itself); call this element $\alpha'$. Then 
the set of all $\alpha'$ is of size $c$ (because of course no two $\alpha'$ are the same), meaning $L_c$ has at least $c$ elements. Thus, by part (d) of this 
exercise (and the Bernstein-Schroder Theorem), it follows that $\overline{\overline{L_c}}=c$. 

\section*{Section 1}

\subsection*{Theorem 1.3}
To me it seems that we need to have the additional assumption that $K$ can be well founded, so that we even know there
is a Mostowski-Shepherdson map over $B$ in the first place. Adding this assumption would not be problematic in 
using Theorem 1.3 to prove Theorem G, as we already know that $L_\alpha$ for all ordinals $\alpha$ is well founded.

Note that $a'=a$ is due to Theorem 5.1 of Chapter 10.

\end{document}
